\documentclass{article}
\usepackage{amsmath}
\begin{document}

\textbf{8. Определяем токи в ветвях исходной схемы методом узловых напряжений}

\begin{align*}
Z_1 &:= 89 + 47i & Z_2 &:= 44 \\
Z_3 &:= -22i & Z_4 &:= 81i \\
Z_5 &:= 55 - 27i & Z_6 &:= -30i \\
Z_7 &:= 63 - 42i \\
E_5 &:= -15.460 + 34.710i & J_6 &:= -3.910 - 5.800i
\end{align*}

\[
A := \begin{pmatrix} 
0 & 0 & 0 & 0 & -1 & 1 & 0 \\ 
1 & 0 & 0 & 0 & 0 & -1 & 0 \\ 
0 & 0 & -1 & 1 & 0 & 0 & 0 \\ 
0 & -1 & 1 & 0 & 0 & 0 & 0 \\ 
0 & 0 & 0 & -1 & 1 & 0 & -1 
\end{pmatrix}
\]

\[
E := \begin{pmatrix} 0 \\ 0 \\ 0 \\ 0 \\ E_5 \\ 0 \\ 0 \end{pmatrix}, \quad 
J_{src} := \begin{pmatrix} 0 \\ 0 \\ 0 \\ 0 \\ 0 \\ J_6 \\ 0 \end{pmatrix}
\]

\[
RD := \text{diag}(Z_1, Z_2, Z_3, Z_4, Z_5, Z_6, Z_7)
\]

\[
G := RD^{-1}
\]

\[
\Phi := (A \cdot G \cdot A^T)^{-1} \cdot (-A \cdot G \cdot E - A \cdot J_{src})
\]

\[
U := A^T \cdot \Phi
\]

\[
IR := G \cdot (U + E) = \begin{pmatrix} 
-0.976 + 0.779i \\ 
-0.132 + 0.863i \\ 
-0.132 + 0.863i \\ 
-0.132 + 0.863i \\ 
-0.976 + 0.779i \\ 
2.938 + 6.582i \\ 
-0.844 - 0.087i 
\end{pmatrix}
\]

\end{document}